% intro_reasonspec.tex — Introduction for the Reasoning Speculation paper
% Loaded by main_reasonspec.tex via % intro_reasonspec.tex — Introduction for the Reasoning Speculation paper
% Loaded by main_reasonspec.tex via % intro_reasonspec.tex — Introduction for the Reasoning Speculation paper
% Loaded by main_reasonspec.tex via % intro_reasonspec.tex — Introduction for the Reasoning Speculation paper
% Loaded by main_reasonspec.tex via \input{sections/intro_reasonspec}

Large language models equipped with chain-of-thought reasoning have achieved remarkable progress on complex tasks by ``thinking longer''---generating extended reasoning traces before committing to an answer.
This paradigm, exemplified by OpenAI o1~\cite{openai2024o1}, QwQ~\cite{qwq2024}, and DeepSeek-R1~\cite{deepseekr1}, frames test-time performance as a scaling problem: allocate more tokens, and accuracy improves.
Yet the gains from extended reasoning are far from uniform across problem instances.
Easy problems are \emph{overthought}: the model wastes hundreds of tokens re-deriving trivial facts or, worse, ``talks itself out of'' a correct intermediate answer through unnecessary self-revision.
Hard problems are \emph{underthought}: a fixed budget truncates reasoning before the model reaches the critical insight.
The fundamental mismatch between uniform compute allocation and heterogeneous problem difficulty leaves substantial accuracy and efficiency gains on the table.

\paragraph{The overthinking problem is pervasive.}
Through systematic experiments across GSM8K with Qwen3-8B, we quantify the extent of this mismatch.
When the reasoning budget is increased from a short to a long allocation, \textbf{XX.X\%} of samples exhibit \emph{accuracy degradation}---they were answered correctly at a lower budget but incorrectly at a higher one.
This is not an isolated edge case; it is a structured phenomenon affecting roughly one-third of all instances.
The ideal solution is clear in principle: allocate compute adaptively, spending little on easy problems and much on hard ones.
An oracle that selects the optimal budget per instance achieves \textbf{XX.X\%} accuracy---far above any fixed budget---demonstrating the enormous headroom for adaptive methods.

\paragraph{Existing adaptive approaches fail in the black-box setting.}
If the gains from adaptive allocation are so large, why has no practical method captured them?
We conducted what is, to our knowledge, the first comprehensive evaluation of black-box adaptive budget controllers for reasoning LLMs.
We tested all major categories of single-path difficulty signals:
(i)~\emph{honest features} extracted from a partial reasoning trace (token utilization, answer presence, lexical coherence), yielding controllers at $-6.5$\,pp below the fixed baseline;
(ii)~\emph{uncertainty-based routing} using answer entropy or self-reported confidence, at $-6.4$\,pp;
and (iii)~\emph{MCTS-style exploration--refinement}, at $-5.6$\,pp.
\textbf{All single-path controllers failed}, often performing \emph{worse} than the na\"ive fixed budget they were designed to improve.
The root cause is informational: a single reasoning trace is fundamentally impoverished as a difficulty signal.
The model's own confidence is poorly calibrated, token-level statistics are noisy, and partial traces reveal little about whether the \emph{final} answer will be correct.
In short, one witness is not enough to judge whether a problem is easy or hard.

\paragraph{Key insight: change the computation structure.}
Rather than trying to extract a weak difficulty signal from a single reasoning path, we propose to \emph{restructure the computation itself} to produce a fundamentally richer signal.
The core idea is simple: instead of running one long chain of thought, first run $K$ short, independent reasoning paths in parallel.
The \emph{cross-path consensus}---the degree to which these $K$ independent ``witnesses'' agree on an answer---provides a difficulty signal that is categorically stronger than anything extractable from a single trace.
Intuitively, if $K$ independent short explorations all converge to the same answer, the problem is almost certainly easy; if they scatter across many answers, it is hard and warrants extended deliberation.
This is analogous in spirit to how speculative decoding~\cite{leviathan2023fast} changed the \emph{structure} of the generation process (small model drafts, large model verifies) without changing the model itself---here, we change the structure of reasoning (parallel exploration, consensus-based routing) without changing the model or its training.
We show in \S\ref{sec:theory} that multi-path agreement provides an information-theoretically optimal difficulty estimator: when each path has independent error probability $\varepsilon$, the consensus signal concentrates exponentially in $K$, yielding a Spearman correlation with true difficulty of $\rho = -0.51$ ($p < 10^{-88}$)---far stronger than any single-path feature we tested.

\paragraph{The \method{} framework.}
We introduce \textbf{\fullmethod{}} (\method{}), a training-free, model-agnostic framework for adaptive test-time compute that operates in three phases:
\begin{enumerate}[nosep,leftmargin=*]
    \item \textbf{Explore.} Given a question $x$, generate $K$ short reasoning paths $\{r_1, \ldots, r_K\}$, each with a small token budget $b_{\text{explore}}$. These paths are independent and can be parallelized.
    \item \textbf{Decide.} Extract the answer from each path and compute a consensus score $c(x) = \max_a \frac{1}{K} \sum_{k=1}^{K} \mathbf{1}[a_k = a]$. Route the problem based on $c(x)$: high consensus $\rightarrow$ accept the majority answer immediately (easy); partial consensus $\rightarrow$ extend the most promising path with additional tokens (medium); low consensus $\rightarrow$ invoke full deliberation with all $K$ paths as context (hard).
    \item \textbf{Resolve.} For medium and hard problems, generate extended reasoning conditioned on the exploration paths, using a budget calibrated to the estimated difficulty tier.
\end{enumerate}
\method{} requires no training, no access to model internals, and no modification to the underlying LLM.
It operates entirely at the API level and is compatible with any reasoning model that supports temperature-based sampling.
The total compute budget is controlled by $K$, $b_{\text{explore}}$, and the per-tier continuation budgets, yielding a flexible accuracy--cost tradeoff.

\paragraph{Contributions.}
Our work makes the following contributions:
\begin{itemize}[nosep,leftmargin=*]
    \item \textbf{A new computation paradigm for adaptive test-time compute.}
    We introduce parallel exploration with consensus-based routing as an alternative to single-path budget allocation.
    Unlike prior adaptive methods that extract features from one reasoning trace, \method{} restructures the computation to produce a multi-path consensus signal that is fundamentally more informative about problem difficulty.

    \item \textbf{Information-theoretic justification.}
    We prove that multi-path consensus dominates any single-path difficulty feature in the black-box setting (\S\ref{sec:theory}).
    The consensus signal concentrates exponentially in $K$, explaining both why single-path controllers fail and why even modest $K$ (e.g., $K=5$) suffices.

    \item \textbf{Comprehensive empirical evaluation.}
    On GSM8K with Qwen3-8B, \method{} achieves \textbf{XX.X\%} accuracy at \textbf{XX\%} of the token cost of a fixed high-budget baseline, representing a \textbf{+X.X\,pp} improvement at matched compute and \textbf{XX\%} token savings at matched accuracy.
    We provide full ablations over $K$, $b_{\text{explore}}$, consensus thresholds, and continuation strategies.

    \item \textbf{First systematic negative results for black-box controllers.}
    We show that \emph{all} major categories of single-path adaptive controllers (honest features, uncertainty routing, MCTS refinement) fail to outperform fixed budgets, with degradations of $5$--$7$\,pp.
    These negative results are themselves a contribution, delimiting what is achievable without structural changes to the computation.

    \item \textbf{Training-free and model-agnostic.}
    \method{} requires no fine-tuning, no reward models, and no access to logits or hidden states.
    It works with any LLM API that supports standard generation and is immediately deployable.
\end{itemize}


Large language models equipped with chain-of-thought reasoning have achieved remarkable progress on complex tasks by ``thinking longer''---generating extended reasoning traces before committing to an answer.
This paradigm, exemplified by OpenAI o1~\cite{openai2024o1}, QwQ~\cite{qwq2024}, and DeepSeek-R1~\cite{deepseekr1}, frames test-time performance as a scaling problem: allocate more tokens, and accuracy improves.
Yet the gains from extended reasoning are far from uniform across problem instances.
Easy problems are \emph{overthought}: the model wastes hundreds of tokens re-deriving trivial facts or, worse, ``talks itself out of'' a correct intermediate answer through unnecessary self-revision.
Hard problems are \emph{underthought}: a fixed budget truncates reasoning before the model reaches the critical insight.
The fundamental mismatch between uniform compute allocation and heterogeneous problem difficulty leaves substantial accuracy and efficiency gains on the table.

\paragraph{The overthinking problem is pervasive.}
Through systematic experiments across GSM8K with Qwen3-8B, we quantify the extent of this mismatch.
When the reasoning budget is increased from a short to a long allocation, \textbf{XX.X\%} of samples exhibit \emph{accuracy degradation}---they were answered correctly at a lower budget but incorrectly at a higher one.
This is not an isolated edge case; it is a structured phenomenon affecting roughly one-third of all instances.
The ideal solution is clear in principle: allocate compute adaptively, spending little on easy problems and much on hard ones.
An oracle that selects the optimal budget per instance achieves \textbf{XX.X\%} accuracy---far above any fixed budget---demonstrating the enormous headroom for adaptive methods.

\paragraph{Existing adaptive approaches fail in the black-box setting.}
If the gains from adaptive allocation are so large, why has no practical method captured them?
We conducted what is, to our knowledge, the first comprehensive evaluation of black-box adaptive budget controllers for reasoning LLMs.
We tested all major categories of single-path difficulty signals:
(i)~\emph{honest features} extracted from a partial reasoning trace (token utilization, answer presence, lexical coherence), yielding controllers at $-6.5$\,pp below the fixed baseline;
(ii)~\emph{uncertainty-based routing} using answer entropy or self-reported confidence, at $-6.4$\,pp;
and (iii)~\emph{MCTS-style exploration--refinement}, at $-5.6$\,pp.
\textbf{All single-path controllers failed}, often performing \emph{worse} than the na\"ive fixed budget they were designed to improve.
The root cause is informational: a single reasoning trace is fundamentally impoverished as a difficulty signal.
The model's own confidence is poorly calibrated, token-level statistics are noisy, and partial traces reveal little about whether the \emph{final} answer will be correct.
In short, one witness is not enough to judge whether a problem is easy or hard.

\paragraph{Key insight: change the computation structure.}
Rather than trying to extract a weak difficulty signal from a single reasoning path, we propose to \emph{restructure the computation itself} to produce a fundamentally richer signal.
The core idea is simple: instead of running one long chain of thought, first run $K$ short, independent reasoning paths in parallel.
The \emph{cross-path consensus}---the degree to which these $K$ independent ``witnesses'' agree on an answer---provides a difficulty signal that is categorically stronger than anything extractable from a single trace.
Intuitively, if $K$ independent short explorations all converge to the same answer, the problem is almost certainly easy; if they scatter across many answers, it is hard and warrants extended deliberation.
This is analogous in spirit to how speculative decoding~\cite{leviathan2023fast} changed the \emph{structure} of the generation process (small model drafts, large model verifies) without changing the model itself---here, we change the structure of reasoning (parallel exploration, consensus-based routing) without changing the model or its training.
We show in \S\ref{sec:theory} that multi-path agreement provides an information-theoretically optimal difficulty estimator: when each path has independent error probability $\varepsilon$, the consensus signal concentrates exponentially in $K$, yielding a Spearman correlation with true difficulty of $\rho = -0.51$ ($p < 10^{-88}$)---far stronger than any single-path feature we tested.

\paragraph{The \method{} framework.}
We introduce \textbf{\fullmethod{}} (\method{}), a training-free, model-agnostic framework for adaptive test-time compute that operates in three phases:
\begin{enumerate}[nosep,leftmargin=*]
    \item \textbf{Explore.} Given a question $x$, generate $K$ short reasoning paths $\{r_1, \ldots, r_K\}$, each with a small token budget $b_{\text{explore}}$. These paths are independent and can be parallelized.
    \item \textbf{Decide.} Extract the answer from each path and compute a consensus score $c(x) = \max_a \frac{1}{K} \sum_{k=1}^{K} \mathbf{1}[a_k = a]$. Route the problem based on $c(x)$: high consensus $\rightarrow$ accept the majority answer immediately (easy); partial consensus $\rightarrow$ extend the most promising path with additional tokens (medium); low consensus $\rightarrow$ invoke full deliberation with all $K$ paths as context (hard).
    \item \textbf{Resolve.} For medium and hard problems, generate extended reasoning conditioned on the exploration paths, using a budget calibrated to the estimated difficulty tier.
\end{enumerate}
\method{} requires no training, no access to model internals, and no modification to the underlying LLM.
It operates entirely at the API level and is compatible with any reasoning model that supports temperature-based sampling.
The total compute budget is controlled by $K$, $b_{\text{explore}}$, and the per-tier continuation budgets, yielding a flexible accuracy--cost tradeoff.

\paragraph{Contributions.}
Our work makes the following contributions:
\begin{itemize}[nosep,leftmargin=*]
    \item \textbf{A new computation paradigm for adaptive test-time compute.}
    We introduce parallel exploration with consensus-based routing as an alternative to single-path budget allocation.
    Unlike prior adaptive methods that extract features from one reasoning trace, \method{} restructures the computation to produce a multi-path consensus signal that is fundamentally more informative about problem difficulty.

    \item \textbf{Information-theoretic justification.}
    We prove that multi-path consensus dominates any single-path difficulty feature in the black-box setting (\S\ref{sec:theory}).
    The consensus signal concentrates exponentially in $K$, explaining both why single-path controllers fail and why even modest $K$ (e.g., $K=5$) suffices.

    \item \textbf{Comprehensive empirical evaluation.}
    On GSM8K with Qwen3-8B, \method{} achieves \textbf{XX.X\%} accuracy at \textbf{XX\%} of the token cost of a fixed high-budget baseline, representing a \textbf{+X.X\,pp} improvement at matched compute and \textbf{XX\%} token savings at matched accuracy.
    We provide full ablations over $K$, $b_{\text{explore}}$, consensus thresholds, and continuation strategies.

    \item \textbf{First systematic negative results for black-box controllers.}
    We show that \emph{all} major categories of single-path adaptive controllers (honest features, uncertainty routing, MCTS refinement) fail to outperform fixed budgets, with degradations of $5$--$7$\,pp.
    These negative results are themselves a contribution, delimiting what is achievable without structural changes to the computation.

    \item \textbf{Training-free and model-agnostic.}
    \method{} requires no fine-tuning, no reward models, and no access to logits or hidden states.
    It works with any LLM API that supports standard generation and is immediately deployable.
\end{itemize}


Large language models equipped with chain-of-thought reasoning have achieved remarkable progress on complex tasks by ``thinking longer''---generating extended reasoning traces before committing to an answer.
This paradigm, exemplified by OpenAI o1~\cite{openai2024o1}, QwQ~\cite{qwq2024}, and DeepSeek-R1~\cite{deepseekr1}, frames test-time performance as a scaling problem: allocate more tokens, and accuracy improves.
Yet the gains from extended reasoning are far from uniform across problem instances.
Easy problems are \emph{overthought}: the model wastes hundreds of tokens re-deriving trivial facts or, worse, ``talks itself out of'' a correct intermediate answer through unnecessary self-revision.
Hard problems are \emph{underthought}: a fixed budget truncates reasoning before the model reaches the critical insight.
The fundamental mismatch between uniform compute allocation and heterogeneous problem difficulty leaves substantial accuracy and efficiency gains on the table.

\paragraph{The overthinking problem is pervasive.}
Through systematic experiments across GSM8K with Qwen3-8B, we quantify the extent of this mismatch.
When the reasoning budget is increased from a short to a long allocation, \textbf{XX.X\%} of samples exhibit \emph{accuracy degradation}---they were answered correctly at a lower budget but incorrectly at a higher one.
This is not an isolated edge case; it is a structured phenomenon affecting roughly one-third of all instances.
The ideal solution is clear in principle: allocate compute adaptively, spending little on easy problems and much on hard ones.
An oracle that selects the optimal budget per instance achieves \textbf{XX.X\%} accuracy---far above any fixed budget---demonstrating the enormous headroom for adaptive methods.

\paragraph{Existing adaptive approaches fail in the black-box setting.}
If the gains from adaptive allocation are so large, why has no practical method captured them?
We conducted what is, to our knowledge, the first comprehensive evaluation of black-box adaptive budget controllers for reasoning LLMs.
We tested all major categories of single-path difficulty signals:
(i)~\emph{honest features} extracted from a partial reasoning trace (token utilization, answer presence, lexical coherence), yielding controllers at $-6.5$\,pp below the fixed baseline;
(ii)~\emph{uncertainty-based routing} using answer entropy or self-reported confidence, at $-6.4$\,pp;
and (iii)~\emph{MCTS-style exploration--refinement}, at $-5.6$\,pp.
\textbf{All single-path controllers failed}, often performing \emph{worse} than the na\"ive fixed budget they were designed to improve.
The root cause is informational: a single reasoning trace is fundamentally impoverished as a difficulty signal.
The model's own confidence is poorly calibrated, token-level statistics are noisy, and partial traces reveal little about whether the \emph{final} answer will be correct.
In short, one witness is not enough to judge whether a problem is easy or hard.

\paragraph{Key insight: change the computation structure.}
Rather than trying to extract a weak difficulty signal from a single reasoning path, we propose to \emph{restructure the computation itself} to produce a fundamentally richer signal.
The core idea is simple: instead of running one long chain of thought, first run $K$ short, independent reasoning paths in parallel.
The \emph{cross-path consensus}---the degree to which these $K$ independent ``witnesses'' agree on an answer---provides a difficulty signal that is categorically stronger than anything extractable from a single trace.
Intuitively, if $K$ independent short explorations all converge to the same answer, the problem is almost certainly easy; if they scatter across many answers, it is hard and warrants extended deliberation.
This is analogous in spirit to how speculative decoding~\cite{leviathan2023fast} changed the \emph{structure} of the generation process (small model drafts, large model verifies) without changing the model itself---here, we change the structure of reasoning (parallel exploration, consensus-based routing) without changing the model or its training.
We show in \S\ref{sec:theory} that multi-path agreement provides an information-theoretically optimal difficulty estimator: when each path has independent error probability $\varepsilon$, the consensus signal concentrates exponentially in $K$, yielding a Spearman correlation with true difficulty of $\rho = -0.51$ ($p < 10^{-88}$)---far stronger than any single-path feature we tested.

\paragraph{The \method{} framework.}
We introduce \textbf{\fullmethod{}} (\method{}), a training-free, model-agnostic framework for adaptive test-time compute that operates in three phases:
\begin{enumerate}[nosep,leftmargin=*]
    \item \textbf{Explore.} Given a question $x$, generate $K$ short reasoning paths $\{r_1, \ldots, r_K\}$, each with a small token budget $b_{\text{explore}}$. These paths are independent and can be parallelized.
    \item \textbf{Decide.} Extract the answer from each path and compute a consensus score $c(x) = \max_a \frac{1}{K} \sum_{k=1}^{K} \mathbf{1}[a_k = a]$. Route the problem based on $c(x)$: high consensus $\rightarrow$ accept the majority answer immediately (easy); partial consensus $\rightarrow$ extend the most promising path with additional tokens (medium); low consensus $\rightarrow$ invoke full deliberation with all $K$ paths as context (hard).
    \item \textbf{Resolve.} For medium and hard problems, generate extended reasoning conditioned on the exploration paths, using a budget calibrated to the estimated difficulty tier.
\end{enumerate}
\method{} requires no training, no access to model internals, and no modification to the underlying LLM.
It operates entirely at the API level and is compatible with any reasoning model that supports temperature-based sampling.
The total compute budget is controlled by $K$, $b_{\text{explore}}$, and the per-tier continuation budgets, yielding a flexible accuracy--cost tradeoff.

\paragraph{Contributions.}
Our work makes the following contributions:
\begin{itemize}[nosep,leftmargin=*]
    \item \textbf{A new computation paradigm for adaptive test-time compute.}
    We introduce parallel exploration with consensus-based routing as an alternative to single-path budget allocation.
    Unlike prior adaptive methods that extract features from one reasoning trace, \method{} restructures the computation to produce a multi-path consensus signal that is fundamentally more informative about problem difficulty.

    \item \textbf{Information-theoretic justification.}
    We prove that multi-path consensus dominates any single-path difficulty feature in the black-box setting (\S\ref{sec:theory}).
    The consensus signal concentrates exponentially in $K$, explaining both why single-path controllers fail and why even modest $K$ (e.g., $K=5$) suffices.

    \item \textbf{Comprehensive empirical evaluation.}
    On GSM8K with Qwen3-8B, \method{} achieves \textbf{XX.X\%} accuracy at \textbf{XX\%} of the token cost of a fixed high-budget baseline, representing a \textbf{+X.X\,pp} improvement at matched compute and \textbf{XX\%} token savings at matched accuracy.
    We provide full ablations over $K$, $b_{\text{explore}}$, consensus thresholds, and continuation strategies.

    \item \textbf{First systematic negative results for black-box controllers.}
    We show that \emph{all} major categories of single-path adaptive controllers (honest features, uncertainty routing, MCTS refinement) fail to outperform fixed budgets, with degradations of $5$--$7$\,pp.
    These negative results are themselves a contribution, delimiting what is achievable without structural changes to the computation.

    \item \textbf{Training-free and model-agnostic.}
    \method{} requires no fine-tuning, no reward models, and no access to logits or hidden states.
    It works with any LLM API that supports standard generation and is immediately deployable.
\end{itemize}


Large language models equipped with chain-of-thought reasoning have achieved remarkable progress on complex tasks by ``thinking longer''---generating extended reasoning traces before committing to an answer.
This paradigm, exemplified by OpenAI o1~\cite{openai2024o1}, QwQ~\cite{qwq2024}, and DeepSeek-R1~\cite{deepseekr1}, frames test-time performance as a scaling problem: allocate more tokens, and accuracy improves.
Yet the gains from extended reasoning are far from uniform across problem instances.
Easy problems are \emph{overthought}: the model wastes hundreds of tokens re-deriving trivial facts or, worse, ``talks itself out of'' a correct intermediate answer through unnecessary self-revision.
Hard problems are \emph{underthought}: a fixed budget truncates reasoning before the model reaches the critical insight.
The fundamental mismatch between uniform compute allocation and heterogeneous problem difficulty leaves substantial accuracy and efficiency gains on the table.

\paragraph{The overthinking problem is pervasive.}
Through systematic experiments across GSM8K with Qwen3-8B, we quantify the extent of this mismatch.
When the reasoning budget is increased from a short to a long allocation, \textbf{XX.X\%} of samples exhibit \emph{accuracy degradation}---they were answered correctly at a lower budget but incorrectly at a higher one.
This is not an isolated edge case; it is a structured phenomenon affecting roughly one-third of all instances.
The ideal solution is clear in principle: allocate compute adaptively, spending little on easy problems and much on hard ones.
An oracle that selects the optimal budget per instance achieves \textbf{XX.X\%} accuracy---far above any fixed budget---demonstrating the enormous headroom for adaptive methods.

\paragraph{Existing adaptive approaches fail in the black-box setting.}
If the gains from adaptive allocation are so large, why has no practical method captured them?
We conducted what is, to our knowledge, the first comprehensive evaluation of black-box adaptive budget controllers for reasoning LLMs.
We tested all major categories of single-path difficulty signals:
(i)~\emph{honest features} extracted from a partial reasoning trace (token utilization, answer presence, lexical coherence), yielding controllers at $-6.5$\,pp below the fixed baseline;
(ii)~\emph{uncertainty-based routing} using answer entropy or self-reported confidence, at $-6.4$\,pp;
and (iii)~\emph{MCTS-style exploration--refinement}, at $-5.6$\,pp.
\textbf{All single-path controllers failed}, often performing \emph{worse} than the na\"ive fixed budget they were designed to improve.
The root cause is informational: a single reasoning trace is fundamentally impoverished as a difficulty signal.
The model's own confidence is poorly calibrated, token-level statistics are noisy, and partial traces reveal little about whether the \emph{final} answer will be correct.
In short, one witness is not enough to judge whether a problem is easy or hard.

\paragraph{Key insight: change the computation structure.}
Rather than trying to extract a weak difficulty signal from a single reasoning path, we propose to \emph{restructure the computation itself} to produce a fundamentally richer signal.
The core idea is simple: instead of running one long chain of thought, first run $K$ short, independent reasoning paths in parallel.
The \emph{cross-path consensus}---the degree to which these $K$ independent ``witnesses'' agree on an answer---provides a difficulty signal that is categorically stronger than anything extractable from a single trace.
Intuitively, if $K$ independent short explorations all converge to the same answer, the problem is almost certainly easy; if they scatter across many answers, it is hard and warrants extended deliberation.
This is analogous in spirit to how speculative decoding~\cite{leviathan2023fast} changed the \emph{structure} of the generation process (small model drafts, large model verifies) without changing the model itself---here, we change the structure of reasoning (parallel exploration, consensus-based routing) without changing the model or its training.
We show in \S\ref{sec:theory} that multi-path agreement provides an information-theoretically optimal difficulty estimator: when each path has independent error probability $\varepsilon$, the consensus signal concentrates exponentially in $K$, yielding a Spearman correlation with true difficulty of $\rho = -0.51$ ($p < 10^{-88}$)---far stronger than any single-path feature we tested.

\paragraph{The \method{} framework.}
We introduce \textbf{\fullmethod{}} (\method{}), a training-free, model-agnostic framework for adaptive test-time compute that operates in three phases:
\begin{enumerate}[nosep,leftmargin=*]
    \item \textbf{Explore.} Given a question $x$, generate $K$ short reasoning paths $\{r_1, \ldots, r_K\}$, each with a small token budget $b_{\text{explore}}$. These paths are independent and can be parallelized.
    \item \textbf{Decide.} Extract the answer from each path and compute a consensus score $c(x) = \max_a \frac{1}{K} \sum_{k=1}^{K} \mathbf{1}[a_k = a]$. Route the problem based on $c(x)$: high consensus $\rightarrow$ accept the majority answer immediately (easy); partial consensus $\rightarrow$ extend the most promising path with additional tokens (medium); low consensus $\rightarrow$ invoke full deliberation with all $K$ paths as context (hard).
    \item \textbf{Resolve.} For medium and hard problems, generate extended reasoning conditioned on the exploration paths, using a budget calibrated to the estimated difficulty tier.
\end{enumerate}
\method{} requires no training, no access to model internals, and no modification to the underlying LLM.
It operates entirely at the API level and is compatible with any reasoning model that supports temperature-based sampling.
The total compute budget is controlled by $K$, $b_{\text{explore}}$, and the per-tier continuation budgets, yielding a flexible accuracy--cost tradeoff.

\paragraph{Contributions.}
Our work makes the following contributions:
\begin{itemize}[nosep,leftmargin=*]
    \item \textbf{A new computation paradigm for adaptive test-time compute.}
    We introduce parallel exploration with consensus-based routing as an alternative to single-path budget allocation.
    Unlike prior adaptive methods that extract features from one reasoning trace, \method{} restructures the computation to produce a multi-path consensus signal that is fundamentally more informative about problem difficulty.

    \item \textbf{Information-theoretic justification.}
    We prove that multi-path consensus dominates any single-path difficulty feature in the black-box setting (\S\ref{sec:theory}).
    The consensus signal concentrates exponentially in $K$, explaining both why single-path controllers fail and why even modest $K$ (e.g., $K=5$) suffices.

    \item \textbf{Comprehensive empirical evaluation.}
    On GSM8K with Qwen3-8B, \method{} achieves \textbf{XX.X\%} accuracy at \textbf{XX\%} of the token cost of a fixed high-budget baseline, representing a \textbf{+X.X\,pp} improvement at matched compute and \textbf{XX\%} token savings at matched accuracy.
    We provide full ablations over $K$, $b_{\text{explore}}$, consensus thresholds, and continuation strategies.

    \item \textbf{First systematic negative results for black-box controllers.}
    We show that \emph{all} major categories of single-path adaptive controllers (honest features, uncertainty routing, MCTS refinement) fail to outperform fixed budgets, with degradations of $5$--$7$\,pp.
    These negative results are themselves a contribution, delimiting what is achievable without structural changes to the computation.

    \item \textbf{Training-free and model-agnostic.}
    \method{} requires no fine-tuning, no reward models, and no access to logits or hidden states.
    It works with any LLM API that supports standard generation and is immediately deployable.
\end{itemize}
